\documentclass[12pt,a4paper]{amsart}

\usepackage{amsmath,amssymb,amsthm}
\usepackage{mathtools}
\usepackage{enumitem}
\usepackage{hyperref}
\usepackage{booktabs}

%%%─── Theorem environments
\newtheorem{theorem}{Theorem}[section]
\newtheorem{lemma}[theorem]{Lemma}
\newtheorem{proposition}[theorem]{Proposition}
\newtheorem{corollary}[theorem]{Corollary}
\newtheorem{conjecture}[theorem]{Conjecture}
\newtheorem{hypothesis}[theorem]{Hypothesis}
\theoremstyle{definition}
\newtheorem{definition}[theorem]{Definition}
\newtheorem{assumption}[theorem]{Assumption}
\newtheorem{problem}[theorem]{Open Problem}
\theoremstyle{remark}
\newtheorem{remark}[theorem]{Remark}
\newtheorem*{remark*}{Remark}

%%%─── Macros
\newcommand{\norm}[1]{\|#1\|}
\newcommand{\ip}[2]{\langle #1,\, #2\rangle}
\newcommand{\Hlim}{\mathcal{H}_{\lim}}
\newcommand{\Hc}{\mathcal{H}_c}
\newcommand{\Phiinf}[1]{\Phi_{#1}^{(\infty)}}
\newcommand{\Phic}[2]{\Phi_{#2}^{(#1)}}
\newcommand{\lambdainf}[1]{\lambda_{#1}^{(\infty)}}
\newcommand{\lambdac}[2]{\lambda_{#2}^{(#1)}}
\newcommand{\PW}{PW}
\newcommand{\spc}[1]{\sigma(#1)}
\newcommand{\Kspace}{\mathcal{K}}

\subjclass[2020]{47A10, 47B25, 42C05, 46B15, 46C05}
\keywords{Prolate spheroidal wave functions, completeness,
  orthonormal basis, strong operator convergence,
  Paley--Wiener space, basis stability, structural classification,
  Hilbert--P\'{o}lya programme}

\begin{document}

\title[Completeness of the Limiting Slepian Eigenbasis]{%
  Completeness of the Limiting Slepian Eigenbasis in $L^2(\mathbb{R})$:
  A Classification of Stability Mechanisms}

\author{[Author]}
\address{[Address]}
\email{[Email]}

\begin{abstract}
Let $\Hlim$ denote the bounded self-adjoint SOT-limit operator
of Paper~X and $\{\Phiinf{n}\}_{n \ge 0}$ the orthonormal
system of limit eigenfunctions of Paper~XI.
Define $\Kspace := \overline{\mathrm{span}}\{\Phiinf{n}\}$.
This paper classifies the mechanisms by which the completeness
problem $\Kspace = L^2(\mathbb{R})$ (Open Problem~7 of Paper~IX)
could be resolved.

\medskip
For finite $c$, the classical PSWFs $\{\Phic{c}{n}\}_{n \ge 0}$
form a complete orthonormal basis of $\PW_c^2$ (Slepian--Pollak
\cite{Slepian1961}).
The central difficulty is that this basis property is not
automatically stable under the SOT-limit $c \to \infty$:
SOT-convergence controls neither spanning structure nor
completeness of the limit system.
Completeness of $\{\Phiinf{n}\}$ is therefore a
\emph{stability problem for PSWF spanning structures under
the SOT-limit}.

\medskip
\textbf{Main results.}
\begin{itemize}
\item \emph{(Theorem~A.)} For any orthonormal system,
  the Parseval identity is logically equivalent to completeness.
  Frame identities are not proof tools.
\item \emph{(Theorem~B.)} Completeness $\Kspace=L^2$ is equivalent
  to the family of local spanning conditions
  $\PW_{c_0}^2 \subseteq \Kspace$ for all $c_0>0$.
  This is a \emph{localisation equivalence}, not a reduction:
  global completeness is rewritten as an
  inductive-limit stability question, with all
  values of $c_0$ playing structurally identical roles.
\item \emph{(Theorem~C.)} Under a compact-resolvent
  hypothesis, completeness follows from the spectral theorem;
  this route is structurally independent of A and B.
\end{itemize}

\medskip
\textbf{Honest assessment.}
This paper does not prove completeness.
It establishes that no frame-based shortcut exists (Route~A),
that completeness is equivalent to a scale-uniform
family of local stability conditions (Route~B),
and that an independent spectral route is available
conditionally (Route~C).
The only currently identified hypothesis strictly weaker than
completeness and directly grounded in the series
is H1 (coefficient stability,
Hypothesis~\ref{hyp:coeff-stability}).
The single-scale formulation of Route~B
(Hypothesis~\ref{hyp:local-span}) is a
\emph{fixed-scale instance} of the full stability problem,
without intrinsic difficulty ordering in $c_0$;
proving it for all $c_0$ is equivalent to completeness.
\end{abstract}

\maketitle

%%═══ 1. INTRODUCTION
\section{Introduction}
\label{sec:intro}

\subsection{The stability problem}

\begin{center}
\fbox{\parbox{0.88\linewidth}{\centering\itshape
For each finite $c$, the PSWFs $\{\Phic{c}{n}\}$ form a
complete orthonormal basis of $\PW_c^2$.
The SOT-limit $c \to \infty$ does not preserve this
basis property automatically.
Completeness of $\{\Phiinf{n}\}$ in $L^2(\mathbb{R})$
is a stability problem:
does the spanning structure of $\{\Phic{c}{n}\}$
survive the limit?}}
\end{center}

\medskip
\noindent
The Slepian concentration operator is
$(\Hc f)(x) = \int_{-\infty}^\infty \frac{\sin c(x-y)}{\pi(x-y)}
f(y)\,dy$.
For each finite $c$, $\{\Phic{c}{n}\}$ is a complete
orthonormal basis of $\PW_c^2$ \cite{Slepian1961}.
Paper~X establishes $\Hlim := \mathrm{SOT\text{-}lim}\,\Hc$
as a bounded self-adjoint operator, $\norm{\Hlim} \le 1$,
$\spc{\Hlim} \subseteq [0,1]$.
Paper~XI constructs $\{\Phiinf{n}\}$ with
\begin{equation}\label{eq:eigen}
  \Hlim\Phiinf{n} = \lambdainf{n}\Phiinf{n},\quad
  \ip{\Phiinf{m}}{\Phiinf{n}} = \delta_{mn},\quad
  \lambdainf{n} \in (0,1).
\end{equation}
The question is
\begin{equation}\label{eq:main}
  \Kspace := \overline{\mathrm{span}}\{\Phiinf{n}\}
  \stackrel{?}{=} L^2(\mathbb{R}).
\end{equation}

\subsection{Why SOT-convergence is not enough}
\label{sub:sot-not-enough}

\begin{itemize}
\item SOT-convergence does not imply
  norm-resolvent convergence \cite[Ex.~VIII.7.14]{ReedSimon1975}.
\item Eigenvalues can dissolve into continuous spectrum
  under SOT-limits \cite[Ex.~VIII.3.4]{ReedSimon1975}.
\item A uniformly bounded family of complete orthonormal
  bases can SOT-converge to a system that is not complete
  (Remark~\ref{rem:sot-stability-failure}).
\end{itemize}

\begin{remark}[SOT-limit can destroy completeness]
\label{rem:sot-stability-failure}
Let $\{e_n\}$ be the standard ONB of $\ell^2$.
For $c < \infty$ define the ONS
$f_n^{(c)} := \cos(c^{-1})e_n + \sin(c^{-1})e_{n+1}$.
Then $f_n^{(c)} \xrightarrow{c\to\infty} e_n$ strongly,
but $\{f_n^{(c)}\}$ is not complete in $\ell^2$ for
any finite $c$ (it misses $e_0 - \tan(c^{-1})e_1$
up to normalisation).
Such behaviour is classical in operator perturbation
theory: a complete ONS can converge SOT to an incomplete
limit (cf.~\cite[Ex.~VIII.3.4]{ReedSimon1975} and
perturbations of ONB projections in
\cite[\S I.4]{Kato1966}).
In the PSWF setting, the finite-$c$ systems are
complete in $\PW_c^2$, not in $L^2$;
the passage $c\to\infty$ changes both the system
and the ambient space, making the stability question
non-trivial.
\end{remark}

\subsection{Honest assessment of this paper}
\label{sub:honest}

\begin{center}
\fbox{\parbox{0.88\linewidth}{
\begin{itemize}
\item Route~B (Theorem~\ref{thm:H2-equiv}) establishes
  an \emph{equivalence}, not a reduction.
  The localisation into $\PW_{c_0}^2$-conditions
  does not simplify the problem;
  it changes the language to one where the
  analytic tools of Papers~VI--XII apply directly.
\item The parameter $c_0$ in the local spanning
  condition plays no intrinsic ordering role:
  proving $\PW_{c_0}^2\subseteq\Kspace$ for one
  $c_0$ is a fixed-scale instance of the stability
  problem, not a canonical intermediate step.
  There is no natural difficulty hierarchy in $c_0$.
\item The only currently identified hypothesis strictly
  weaker than completeness and grounded in the series
  is~H1.
\item This paper is a barrier and classification paper,
  not a progress paper.
\end{itemize}}}
\end{center}

\subsection{Dependency graph}
\label{sub:deps}
\[
  \underbrace{\text{H1}}_{\text{Hyp.~\ref{hyp:coeff-stability}}}
  \;\text{strictly weaker than completeness},
\]
\[
  \underbrace{\text{local H2}_{c_0}}_{\text{Hyp.~\ref{hyp:local-span}}}
  \;\text{strictly weaker (fixed-scale instance, no }c_0\text{-ordering)},
\]
\[
  \underbrace{\text{H2 (all }c_0\text{)}}_{\text{Thm.~\ref{thm:H2-equiv}}}
  \;\Longleftrightarrow\;
  \Kspace = L^2
  \;\Longleftrightarrow\;
  \text{Route~A (Thm.~\ref{thm:routeA})},
\]
\[
  \underbrace{\text{compact resolvent}}_{\text{Hyp.~\ref{hyp:compact-resolvent}}}
  \xRightarrow{\text{Prop.~\ref{prop:compact-route}}}
  \Kspace = L^2
  \quad\text{(structurally independent)}.
\]

\subsection{What this paper does not claim}
\begin{itemize}
\item We do not prove $\Kspace = L^2(\mathbb{R})$.
\item We do not prove local H2$_{c_0}$ for any $c_0$.
\item We do not claim $\spc{\Hlim}$ is discrete.
\item We do not claim local H2$_{c_0}$ is easier than
  completeness in any intrinsic sense.
\end{itemize}

\subsection{Role in the series}
\begin{center}
\begin{tabular}{lll}
\toprule
Paper & Content & Status \\
\midrule
X  & SOT-limit; Mosco; Friedrichs & Final \\
XI & ONS $\{\Phiinf{n}\}$; orthonormality & Final \\
XII & Localisation; obstruction set & Final \\
\textbf{XIII} & \textbf{Classification of stability mechanisms}
    & \textbf{This paper} \\
\bottomrule
\end{tabular}
\end{center}

%%═══ 2. INHERITED DATA
\section{Inherited Data}
\label{sec:input}

\begin{lemma}[Classical spanning at finite $c$]
\label{lem:finite-c-basis}
For each finite $c > 0$,
$\{\Phic{c}{n}\}_{n \ge 0}$ is a complete orthonormal
basis of $\PW_c^2$ \cite{Slepian1961,Osipov2013}.
In particular,
$\sum_{n \ge 0}|\ip{g}{\Phic{c}{n}}|^2 = \norm{g}^2$
for all $g \in \PW_c^2$.
\end{lemma}

\begin{lemma}[Series results used]
\label{lem:inherited}
\begin{enumerate}[\rm(A)]
\item\label{A:sa}
  \emph{(Paper~X.)} $\Hlim$ bounded self-adjoint,
  $\norm{\Hlim}\le 1$, $\spc{\Hlim}\subseteq[0,1]$.
\item\label{A:ONS}
  \emph{(Paper~XI.)} $\{\Phiinf{n}\}$ satisfies
  \eqref{eq:eigen} and is orthonormal.
\item\label{A:nonvanish}
  \emph{(Paper~VIII.)} For $n\le\kappa c$:
  $|\widehat{\Phic{c}{n}}|\ge c_1 c^{-1/2}$ on $E_c$,
  $|E_c|\ge c_* c^{-1/2}$.
\item\label{A:precompact}
  \emph{(Paper~IX, Lem.~3.1.)} $\{\Phic{c}{n}\}_{c>0}$
  precompact in $L^2$.
\item\label{A:eval-rate}
  \emph{(Paper~VIII.)} $|\lambdac{c}{n}-\lambdainf{n}|
  \le C_\kappa c^{-1/4}$ for $n\le\kappa c$.
\item\label{A:evec-conv}
  \emph{(Paper~XII, cond.~Gap-S+SOT-Faster.)}
  $\norm{\Phic{c}{n}-\Phiinf{n}}\to 0$ for each fixed $n$.
\end{enumerate}
\end{lemma}

\begin{remark}[The stability gap]
\label{rem:stability-gap}
Lemma~\ref{lem:finite-c-basis} gives the complete
Parseval identity on each $\PW_c^2$ for finite $c$.
Lemma~\ref{lem:inherited}\ref{A:evec-conv} gives
pointwise convergence $\Phic{c}{n}\to\Phiinf{n}$
(conditionally).
The gap is:
\emph{pointwise convergence of basis vectors does
not imply preservation of the basis property in the
limit space.}
This is the stability problem at the heart of Paper~XIII.
\end{remark}

%%══─ 3. INTERFACE LEMMA
\section{Interface Lemma}
\label{sec:interface}

\begin{lemma}[Orthogonality characterisation]
\label{lem:orth-char}
$\Kspace = L^2(\mathbb{R})\iff
[\ip{f}{\Phiinf{n}}=0\;\forall n\Rightarrow f=0]$.
\end{lemma}
\begin{proof}\cite[Thm.~3.3]{Conway1990}.\end{proof}

%%═══ 4. ROUTE A
\section{Route~A: Parseval is Tautologically Equivalent
  to Completeness}
\label{sec:routeA}

\begin{definition}[Gram operator]
\label{def:gram}
$Gf := \sum_{n\ge 0}\ip{f}{\Phiinf{n}}\Phiinf{n}$.
Bessel's inequality gives
$\norm{S_Nf-S_Mf}^2 = \sum_{n=M+1}^N
|\ip{f}{\Phiinf{n}}|^2\to 0$,
so the series converges in $L^2$-norm.
$G$ is self-adjoint ($G^*=G$) and idempotent
($G^2=G$) by orthonormality alone,
hence the orthogonal projection onto $\Kspace$.
(This uses only orthonormality, not completeness.)
\end{definition}

\begin{theorem}[Route~A]
\label{thm:routeA}
For the orthonormal system $\{\Phiinf{n}\}$:
\[
  \Kspace = L^2(\mathbb{R})
  \iff G = \mathrm{id}
  \iff \textstyle\sum_n|\ip{f}{\Phiinf{n}}|^2 = \norm{f}^2\ \forall f
  \iff \exists A{>}0:\,
    \textstyle\sum_n|\ip{f}{\Phiinf{n}}|^2 \ge A\norm{f}^2\ \forall f.
\]
\end{theorem}
\begin{proof}
(i)$\Rightarrow$(ii): if $\Kspace=L^2$ then $G=\mathrm{id}$.
(ii)$\Rightarrow$(iii): $\norm{Gf}^2=\norm{f}^2$.
(iii)$\Rightarrow$(iv): take $A=1$.
(iv)$\Rightarrow$(i): if $f\perp\Phiinf{n}$ for all $n$
then $0\ge A\norm{f}^2$, so $f=0$;
hence $\Kspace=L^2$ by Lemma~\ref{lem:orth-char}.
\end{proof}

\begin{remark}[Route~A is a closure, not a starting point]
\label{rem:routeA-consequence}
Theorem~\ref{thm:routeA} shows that every
frame-based formulation of completeness is
tautologically equivalent to completeness itself.
The correct logical direction is:
prove $\Kspace=L^2$ (by density or spanning arguments),
then deduce Parseval as a consequence.
Attempts to use Parseval or a frame lower bound
\emph{to prove} completeness must implicitly already
know completeness.
This closes off all frame-based proof strategies;
any approach must work at the level of individual
$\PW_c^2$-bands before passing to $L^2$ by density.
\end{remark}

%%═══ 5. ROUTE B
\section{Route~B: Completeness as a Stability Problem
  for PSWF Spanning Structures}
\label{sec:routeB}

\subsection{The localisation equivalence}

\begin{lemma}[Paley--Wiener density]
\label{lem:pw-dense}
$\bigcup_{c>0}\PW_c^2$ is dense in $L^2(\mathbb{R})$.
\end{lemma}
\begin{proof}
Schwartz functions with compact Fourier support lie
in $\bigcup_c\PW_c^2$ and are dense
\cite[Thm.~7.14]{Rudin1991}.\end{proof}

\begin{theorem}[Route~B: localisation equivalence]
\label{thm:H2-equiv}
The following are equivalent:
\begin{enumerate}[\rm(i)]
\item $\Kspace = L^2(\mathbb{R})$;
\item $\PW_{c_0}^2 \subseteq \Kspace$ for every $c_0 > 0$.
\end{enumerate}
\end{theorem}
\begin{proof}
(i)$\Rightarrow$(ii): trivial.
(ii)$\Rightarrow$(i): $\bigcup_{c_0}\PW_{c_0}^2
\subseteq\Kspace$;
since $\Kspace$ is closed and
$\bigcup_{c_0}\PW_{c_0}^2$ is dense
(Lemma~\ref{lem:pw-dense}),
we get $L^2=\overline{\bigcup_{c_0}\PW_{c_0}^2}
\subseteq\Kspace$.
\end{proof}

\begin{remark}[Route~B: equivalence, not reduction]
\label{rem:routeB-equiv}
Theorem~\ref{thm:H2-equiv} is a localisation
\emph{equivalence}: global completeness is precisely
the family of local spanning conditions.
This is not a reduction of difficulty:
the quantification over all $c_0 > 0$ makes the
local formulation exactly as strong as global completeness.

The parameter $c_0$ plays no intrinsic ordering role.
Proving $\PW_{c_0}^2 \subseteq \Kspace$ for one
fixed $c_0$ is a \emph{fixed-scale instance}
of the stability problem, with the same structural
character for every value of $c_0$;
there is no canonical hierarchy of difficulty in $c_0$.

The \emph{conceptual value} of Route~B is not logical
simplification but \emph{language change}:
each individual $\PW_{c_0}^2$ is a concrete
separable Hilbert space with explicit finite-$c_0$
PSWF basis (Lemma~\ref{lem:finite-c-basis}),
and the question becomes a stability question
about that basis under the SOT-limit.
This is the language in which the analytic
tools of Papers~VI--XII (Fourier decay, eigenvector
convergence, nonvanishing) apply most directly.
\end{remark}

\subsection{The fixed-scale instance}

\begin{hypothesis}[Local spanning: fixed scale]
\label{hyp:local-span}
For a single fixed $c_0 > 0$:
$\PW_{c_0}^2 \subseteq \Kspace$.
\end{hypothesis}

\begin{remark}[Status of Hypothesis~\ref{hyp:local-span}]
\label{rem:local-H2-status}
Hypothesis~\ref{hyp:local-span} for one fixed $c_0$
is strictly weaker than completeness
($\Kspace=L^2$ implies it, but one instance does
not imply $\Kspace=L^2$).
However, it is a \emph{fixed-scale instance} of
the stability problem, not a canonical intermediate
step:
all values of $c_0$ have the same structural
character, and there is no intrinsic ordering
of difficulty in $c_0$.
A proof for all $c_0$ would, by
Theorem~\ref{thm:H2-equiv}, yield completeness;
but this uniform quantification is the hard step,
not any individual instance.
The natural approach to each instance is:
start from the finite-$c_0$ basis of
Lemma~\ref{lem:finite-c-basis} and control
the approximation error under $c \to \infty$
via Papers~VI--VIII.
\end{remark}

\subsection{Coefficient stability}

\begin{hypothesis}[H1: coefficient stability]
\label{hyp:coeff-stability}
For each $c_0>0$ and $g\in\PW_{c_0}^2$,
$a_n^{(c)}:=\ip{g}{\Phic{c}{n}}
\to a_n^{(\infty)}:=\ip{g}{\Phiinf{n}}$
in $\ell^2(\mathbb{N})$ as $c\to\infty$.
\end{hypothesis}

\begin{remark}[H1 is the only currently identified strictly
  weaker grounded hypothesis]
\label{rem:H1-status}
H1 concerns $\ell^2$-stability of \emph{coordinates}
under the SOT-limit.
It does not assert reconstruction of $g$ in $L^2$
(that would be H2$_{c_0}$, equivalent to completeness
if true for all $c_0$).
H1 is therefore the only currently identified
hypothesis in this paper that is:
\begin{itemize}
\item strictly weaker than completeness;
\item directly supported by the series:
  pointwise convergence from Paper~XII
  (Lemma~\ref{lem:inherited}\ref{A:evec-conv}),
  non-collapse from Paper~VIII
  (Lemma~\ref{lem:inherited}\ref{A:nonvanish}),
  uniform $\ell^2$-bound from Bessel.
\end{itemize}
The open step is norm convergence in $\ell^2$:
\begin{equation}\label{eq:norm-H1}
  \textstyle\sum_n|a_n^{(c)}|^2
  \to\sum_n|a_n^{(\infty)}|^2.
\end{equation}
This follows from weak $\ell^2$-convergence plus
norm convergence via the Hilbert-space identity
$\norm{x-y}^2=\norm{x}^2-2\,\mathrm{Re}\ip{x}{y}+\norm{y}^2$;
the missing step is \eqref{eq:norm-H1} itself,
for which Papers~VIII+XII provide the natural
starting point.
\end{remark}

\subsection{Conditional completeness}

\begin{theorem}[Conditional completeness via Route~B]
\label{thm:main-B}
\textup{(Conditional on H2$_{c_0}$ for all $c_0>0$.)}
If $\PW_{c_0}^2\subseteq\Kspace$ for all $c_0>0$,
then $\Kspace = L^2(\mathbb{R})$.
\end{theorem}
\begin{proof}
Immediate from Theorem~\ref{thm:H2-equiv}.
\end{proof}

\begin{remark}[Role of H1]
\label{rem:H1-role}
H1 does not appear in the proof of
Theorem~\ref{thm:main-B}.
It is relevant as \emph{necessary condition} for
H2$_{c_0}$: if $g\in\PW_{c_0}^2$ is reconstructed
by $\{\Phiinf{n}\}$, then its coordinates must be
$\ell^2$-stable.
H1 is thus a verifiable partial step
\emph{towards} H2$_{c_0}$, not a proof of it.
\end{remark}

%%═══ 6. ROUTE C
\section{Route~C: Conditional Compactness}
\label{sec:routeC}

\begin{remark}[Route~C is structurally independent]
\label{rem:routeC-indep}
Routes~A and~B operate within the PSWF geometry
($\PW_c^2$, finite-$c$ basis, Papers~VI--XII).
Route~C posits a global compactness hypothesis
that, if true, gives completeness from the
spectral theorem alone.
The two approaches reflect different operator
regimes and are independent alternatives.
\end{remark}

\begin{hypothesis}[Compact resolvent]
\label{hyp:compact-resolvent}
$(\Hlim + I)^{-1}$ is compact on $L^2(\mathbb{R})$.
\end{hypothesis}

\begin{proposition}[Route~C]
\label{prop:compact-route}
\textup{(Conditional on Hypothesis~\ref{hyp:compact-resolvent}.)}
If $(\Hlim+I)^{-1}$ is compact, then $\Hlim$ has
purely discrete spectrum with finite multiplicities
accumulating only at $0$ or $1$,
$\{\Phiinf{n}\}$ is an orthonormal basis,
and $\Kspace = L^2(\mathbb{R})$.
\end{proposition}
\begin{proof}
$T:=(\Hlim+I)^{-1}$ is compact self-adjoint.
By \cite[Thm.~VI.3.2]{ReedSimon1975},
$\sigma(T)$ consists of eigenvalues accumulating at $0$.
Since $\lambdainf{n}>0$, the eigenvalues of $T$ are
$(\lambdainf{n}+1)^{-1}>0$,
so $\spc{\Hlim}$ is purely discrete;
$\spc{\Hlim}\subseteq[0,1]$ forces accumulation
only at $0$ or $1$.
The spectral theorem \cite[Thm.~VIII.6]{ReedSimon1975}
gives the ONB property.
\end{proof}

\begin{remark}[Incompatibility with $\spc{\Hlim}=[0,1]$]
\label{rem:op4}
Compact resolvent and continuous spectrum $[0,1]$
are mutually exclusive.
Paper~XIII does not resolve this tension
(Paper~X, Open Problem~4).
\end{remark}

%%═══ 7. CLASSIFICATION
\section{Classification}
\label{sec:classification}

\begin{theorem}[Classification of completeness mechanisms]
\label{thm:classification}
\hfill
\begin{center}
\renewcommand{\arraystretch}{1.4}
\begin{tabular}{lllp{4.5cm}}
\toprule
Mechanism & Logical type & Strength & Analytic core \\
\midrule
Route~A & Tautological equivalence
  & $=$ completeness
  & No proof route via frames exists \\
Route~B & Localisation equivalence
  & $=$ completeness
  & SOT-stability of $\PW_{c_0}^2$-bases
    (no $c_0$-ordering) \\
Local H2$_{c_0}$ & Fixed-scale instance
  & $<$ completeness
  & Basis stability at fixed scale;
    no canonical hierarchy in $c_0$ \\
H1 & Coefficient stability
  & $<$ completeness
  & Only currently identified
    grounded weaker hypothesis \\
Route~C & Independent structural route
  & independent
  & Requires compact resolvent \\
\bottomrule
\end{tabular}
\end{center}
\end{theorem}

%%═══ 8. OPEN PROBLEMS
\section{Open Problems}
\label{sec:open}

\begin{problem}[H1: coefficient norm convergence]
\label{prob:H1}
\emph{(Strictly weaker than completeness; most accessible.)}
Prove~\eqref{eq:norm-H1}.
Natural tools: Paper~VIII nonvanishing,
Paper~XII eigenvector convergence, Bessel majorant.
\end{problem}

\begin{problem}[Fixed-scale spanning]
\label{prob:local-span}
\emph{(Strictly weaker than completeness;
fixed-scale instance; no intrinsic $c_0$-ordering.)}
For a fixed but arbitrary $c_0 > 0$
(no canonical ordering in $c_0$),
prove $\PW_{c_0}^2\subseteq\Kspace$.
By Theorem~\ref{thm:H2-equiv},
proving this for all $c_0$ yields completeness.
Natural approach: $\{\Phic{c_0}{n}\}$ spans
$\PW_{c_0}^2$ for finite $c_0$; control the limit
via Papers~VI--VIII.
\end{problem}

\begin{problem}[Spectral type of $\Hlim$]
\label{prob:spectral-type}
\emph{(Hardest; independent of Routes~A+B.)}
Is $\spc{\Hlim}$ purely discrete?
Is $\spc{\Hlim}=[0,1]$?
Is $(\Hlim+I)^{-1}$ compact?
The last two are mutually exclusive.
\end{problem}

\begin{problem}[Operator identification]
\label{prob:bridge}
Prove $\Hlim = \overline{H_{\mathrm{spec}}}$
(Paper~IX OP~7; Paper~X Bridge Theorem).
Requires completeness of $\{\Phiinf{n}\}$.
\end{problem}

%%═══ 9. SUMMARY
\section{Summary}
\label{sec:summary}

\medskip
\begin{center}
\renewcommand{\arraystretch}{1.3}
\begin{tabular}{lllp{3.5cm}}
\toprule
Mechanism & Strength & Status & Core \\
\midrule
Route~A & $=$ completeness
  & Closed (tautology)
  & No frame shortcut \\
Route~B & $=$ completeness
  & Open
  & SOT-stability of $\PW_c^2$-bases \\
Local H2$_{c_0}$ & $<$ completeness
  & Open
  & Fixed-scale instance; no $c_0$-hierarchy \\
H1 & $<$ completeness
  & Open
  & Norm conv.\ \eqref{eq:norm-H1} \\
Route~C & independent
  & Open (conditional)
  & Compact resolvent \\
\bottomrule
\end{tabular}
\end{center}

\medskip
\noindent
This paper is a \emph{barrier and classification paper}.
The central conclusion is:
\emph{completeness of $\{\Phiinf{n}\}$ is a stability
problem for PSWF spanning structures under the
SOT-limit, and all tautological and frame-based
approaches are closed off.}
The open core of the problem remains the stability
of individual $\PW_{c_0}^2$-bases under
$c\to\infty$ (Route~B / local H2),
with H1 as the only currently identified grounded
weaker step.

\begin{thebibliography}{99}

\bibitem{PaperXIII-XI}[Author],
\textit{Spectral structure, Paper~XI}, 2026.

\bibitem{PaperXIII-X}[Author],
\textit{Mosco convergence and Friedrichs, Paper~X}, 2026.

\bibitem{PaperXIII-IX}[Author],
\textit{Conditional framework, Paper~IX}, 2026.

\bibitem{PaperXIII-VIII}[Author],
\textit{Non-cancellation and lower bounds, Paper~VIII}, 2026.

\bibitem{PaperXIII-XII}[Author],
\textit{Localisation principle, Paper~XII}, 2026.

\bibitem{Slepian1961}
D.~Slepian, H.~O.~Pollak,
\textit{Prolate spheroidal wave functions~I},
Bell Syst.\ Tech.\ J.\ \textbf{40} (1961), 43--63.

\bibitem{Osipov2013}
A.~Osipov, V.~Rokhlin, H.~Xiao,
\textit{Prolate Spheroidal Wave Functions of Order Zero},
Springer, 2013.

\bibitem{Kato1966}
T.~Kato, \textit{Perturbation Theory for Linear Operators},
Springer, 1966.

\bibitem{ReedSimon1975}
M.~Reed, B.~Simon,
\textit{Methods of Modern Mathematical Physics, Vol.~II},
Academic Press, 1975.

\bibitem{Conway1990}
J.~B.~Conway,
\textit{A Course in Functional Analysis}, 2nd ed.,
Springer, 1990.

\bibitem{Rudin1991}
W.~Rudin, \textit{Functional Analysis}, 2nd ed.,
McGraw-Hill, 1991.

\end{thebibliography}

\end{document}
